% BHU PYQ -- Manually transcribed & error-corrected
% Paper: GLB-603: Hydrogeology, Environmental Geology, Exploration Geology & Computer Applications
% Exam: B.Sc. Semester VI Examination, 2013-14

\documentclass[11pt,a4paper]{article}
\usepackage[a4paper,top=2.5cm,bottom=2.5cm,left=2.8cm,right=2.8cm,headheight=14pt]{geometry}
\usepackage{amsmath,amssymb}
\usepackage[T1]{fontenc}
\usepackage[utf8]{inputenc}
\usepackage{lmodern,microtype,fancyhdr,enumitem,hyperref,lastpage,parskip}

\hypersetup{colorlinks=true,urlcolor=black,linkcolor=black,
  pdftitle={GLB-603 Hydrogeology Environmental Geology Exploration Geology Computer Applications -- BHU B.Sc. Sem VI 2013-14}}
\pagestyle{fancy}\fancyhf{}
\renewcommand{\headrulewidth}{0.4pt}\renewcommand{\footrulewidth}{0.4pt}
\fancyhead[L]{\small   \textit{Banaras Hindu University}}
\fancyhead[R]{\small   \textit{GLB-603: Hydrogeology, Environmental \& Exploration Geology}}
\fancyfoot[C]{\small Page  \thepage\ of \pageref{LastPage}}


\newlist{parts}{enumerate}{2}
\setlist[parts,1]{label=(\alph*),leftmargin=*,itemsep=5pt,topsep=3pt}

\newlist{subparts}{enumerate}{2}
\setlist[subparts,1]{label=(\roman*),leftmargin=*,itemsep=3pt,topsep=2pt}

\newcommand{\pts}[1]{\hfill   \textit{[#1]}}


\begin{document}

\begin{center}
  {\large\bfseries BANARAS HINDU UNIVERSITY}\\[2pt]
  {
\normalsize B.Sc. --- Semester VI Examination, 2013-14}\\[6pt]
  \rule{0.8\linewidth}{0.5pt}\\[6pt]
  {\large\bfseries Paper: GLB-603}\\[2pt]
  {
ormalsizefseries (Hydrogeology, Environmental Geology, Exploration Geology}\\[2pt]
  {
ormalsizefseries and Computer Applications)}\\[4pt]
  {
\normalsize Time: 3 Hours \hfill Full Marks: 70}\\
\end{center}
\rule{\linewidth}{0.4pt}
\smallskip



\noindent   \textit{Note: Attempt   \textbf{five} questions in all, selecting at least   \textbf{two} from each   \textbf{Section-A} and   \textbf{B}.   \textbf{Section-C} is   \textit{compulsory}. All questions carry equal marks.}
\bigskip

\bigskip


\noindent {\large\bfseries SECTION --- A}
\rule{\linewidth}{0.4pt}\smallskip




\noindent   \textbf{Question 1.} Discuss the salient features of Geochemical Method of Mineral Exploration giving suitable examples. \pts{14}

\medskip


\noindent   \textbf{Question 2.} Write a brief note on the Important Features of MS PowerPoint. \pts{14}

\medskip


\noindent   \textbf{Question 3.} What is Hydrologic Cycle? Give a brief account of processes involved in it. Supplement your answer with a suitable diagram. \pts{14}

\bigskip


\noindent {\large\bfseries SECTION --- B}
\rule{\linewidth}{0.4pt}\smallskip




\noindent   \textbf{Question 4.} Describe briefly: \pts{14}

\begin{parts}
  \item Electrical method of exploration
  \item Indicator plants
  \item Economic criteria in mineral exploration
\end{parts}

\medskip


\noindent   \textbf{Question 5.} Answer the following questions: \pts{14}

\begin{parts}
  \item Explain how you would insert a table in a Word document.
  \item Explain the procedure of addition of numbers in MS Excel with example.
  \item How do you export photographic plate in Corel Draw?
  \item Write a short note on the portable storage devices of computer.
\end{parts}

\medskip


\noindent   \textbf{Question 6.} Differentiate between the following: \pts{14}

\begin{parts}
  \item Unconfined and Confined aquifers
  \item Aquitard and aquifuge
  \item Climate and Weather
  \item Fall and Debris Flow
\end{parts}

\medskip


\noindent   \textbf{Question 7.} Write explanatory notes on the following: \pts{14}

\begin{parts}
  \item Seismic waves
  \item Noise pollution
\end{parts}

\bigskip


\noindent {\large\bfseries SECTION --- C}
\rule{\linewidth}{0.4pt}\smallskip




\noindent   \textbf{Question 8.} Answer/Explain the following very briefly (not more than 50 words): \pts{$7  \times 2 = 14$}

\begin{parts}
  \item What are the mutational effects of iron and copper?
  \item What is Reconnaissance Permit (RP)?
  \item Permeability
  \item Reverse osmosis
  \item Flash flood
  \item What is the computer system?
  \item What is the difference between save and save as command in MS Office?
\end{parts}

\vfill\rule{\linewidth}{0.4pt}

\begin{center}\small
\href{https://bhu-exam.vercel.app/}{Click here to visit our website}\\[4pt]
\href{https://me-aryan.vercel.app/}{Click here to visit our contributor}\\[4pt]
\href{https://quashan-me.vercel.app/}{Click here to visit our contributor}
\end{center}
\end{document}
